% \newcommand*{\StylePath}{../styles}
% % 2014.08.05: This preamble block is now moved to each individual main file, because \import is now used instead of \input
% \documentclass[12pt]{article}
% \usepackage{amsmath,amssymb}
% \usepackage{color,graphicx,fancybox,pause}
% \usepackage[utf8]{vntex} % type Vietnamese
% \usepackage{sagetex} % run Sagemath within Latex

\renewcommand{\familydefault}{cmss}
\renewcommand{\floatpagefraction}{1}

%------------------------------------------
%  PDF
%------------------------------------------
\pdfoutput=1
\pdfcompresslevel=0
\pdfimageresolution=600
\pdfpkresolution=600
\pagestyle{empty}
\pdfpageheight=11.1cm
\pdfpagewidth=14.8cm
%------------------------------------------
% Text dimensions
%------------------------------------------
\textheight=12cm
\textwidth=12.2cm
\topmargin-3cm
\oddsidemargin-1.2cm
\parindent0cm
\parskip0ex
%\arraycolsep0.4ex
%------------------------------------------
%  Lengths
%------------------------------------------
\renewcommand{\baselinestretch}{1.2} % These two lines, although not called in .tex file, affect the output
\renewcommand{\arraystretch}{1}
\newcommand{\spa}[1]{\par\vspace*{#1ex}}
\newcommand{\skipthis}[1]{}

%Set standard display distances
%\rennewcommand{\arraystretch}{1.1}

\newcommand{\sdd}{
\setlength{\jot}{1ex}
\setlength{\abovedisplayskip}{1ex}
\setlength{\abovedisplayshortskip}{1ex}
\setlength{\belowdisplayskip}{1ex}
\setlength{\belowdisplayshortskip}{1ex}}

\newlength{\tew}\setlength{\tew}{\textwidth}
\newlength{\pah}\setlength{\pah}{\pdfpageheight}
\newlength{\storea}
\newlength{\storeb}
\newlength{\storec}
\newlength{\stored}
\newlength{\storee}
\newlength{\storef}

%Scale display distance box
\newcommand{\sdb}[2]{
\storea\abovedisplayskip
\storeb\abovedisplayshortskip
\storec\belowdisplayskip
\stored\belowdisplayshortskip
\setlength{\abovedisplayskip}{#1\abovedisplayskip}
\setlength{\abovedisplayshortskip}{#1\abovedisplayshortskip}
\setlength{\belowdisplayskip}{#1\belowdisplayskip}
\setlength{\belowdisplayshortskip}{#1\belowdisplayshortskip}

#2

\setlength{\abovedisplayskip}{\storea}
\setlength{\abovedisplayshortskip}{\storeb}
\setlength{\belowdisplayskip}{\storec}
\setlength{\belowdisplayshortskip}{\stored}
}

%Scale display distance box with different factors top and bottom
\newcommand{\sdbe}[3]{
\storea=\abovedisplayskip
\storeb=\abovedisplayshortskip
\storec=\belowdisplayskip
\stored=\belowdisplayshortskip
\setlength{\abovedisplayskip}{#1\abovedisplayskip}
\setlength{\abovedisplayshortskip}{#1\abovedisplayshortskip}
\setlength{\belowdisplayskip}{#2\belowdisplayskip}
\setlength{\belowdisplayshortskip}{#2\belowdisplayshortskip}

#3

\setlength{\abovedisplayskip}{\storea}
\setlength{\abovedisplayshortskip}{\storeb}
\setlength{\belowdisplayskip}{\storec}
\setlength{\belowdisplayshortskip}{\stored}
}

%Baselineskip box: scale baslineskip
\newcommand{\blb}[2]{
\storee=\baselineskip
\setlength{\baselineskip}{#1\baselineskip}
#2

\setlength{\baselineskip}{\storee}
}

%------------------------------------------
\newcommand{\enu}[1]{ \begin{enumerate} #1 \end{enumerate} }
\newcommand{\ite}[1]{\vspace*{0ex}\begin{list}{$\bullet$}{\leftmargin2.5ex}#1\end{list} } % itemize list without indentation
\newcommand{\itemi}[1]{\begin{itemize}#1\end{itemize} }

\renewcommand{\labelenumi}{\arabic{enumi}.}
\renewcommand{\labelenumii}{(\alph{enumii})}

\newcommand{\aln}[1]{\noindent\begin{align*}#1\end{align*}}
\newcommand{\mun}[1]{\noindent\begin{multline*}#1\end{multline*}}
\newcommand{\eqn}[1]{\spa{-3}$$#1$$\spa{0}}
\newcommand{\eqns}[1]{\spa{-2}$$#1$$\spa{0}}
\newcommand{\gan}[1]{\begin{gather*}#1\end{gather*}}
\newcommand{\ear}[1]{\begin{eqnarray}#1\end{eqnarray}}
\newcommand{\ean}[1]{\begin{eqnarray*}#1\end{eqnarray*}}

\newcommand{\arr}[2]{\begin{array}{#1}#2\end{array}}
\newcommand{\tab}[2]{\begin{tabular}{#1}#2\end{tabular}}
\newcommand{\mas}[2]{\left[\begin{array}{#1}#2\end{array}\right]}
\newcommand{\mat}[2]{\left(\begin{array}{#1}#2\end{array}\right)}
%\newcommand{\mat}[2]{\mas{#1}{#2}}
\newcommand{\matl}[4]{\left#1\begin{array}{#2}#3\end{array}\right#4}
\newcommand{\dmas}[2]{
\mas{ccc}{#1&&0\\[-1.5ex]&\ddots&\\[-1.5ex]0&&#2}}
\newcommand{\quo}[1]{\noindent\vspace*{-1ex}\begin{quote}#1\end{quote}}
\newcommand{\cen}[1]{\begin{center}#1\end{center}}
\newcommand{\epro}{ \mbox{\rule{2mm}{2mm}} }

%-----------------------------------------------------------------------
%title frontpage CTUT and STW
\newcommand{\titfstw}[2]{%\vspace*{-1ex}
{\hspace*{-0.0\tew}{\bf\Large\color{black}{{#1}}}\\[-1.6ex] }\\[0.70\pah]
\fboxsep=0.3ex
\hspace*{-0.12\tew}
\colorbox{mag}{
\begin{minipage}{1.25\tew}
\tiny  \hfill \phantom{\arabic{page}/\pageref{lastpage}}\hspace*{1cm}
\end{minipage}}
\hspace*{-0.12\tew}
\begin{minipage}{1.25\tew}
\vspace*{-0.7ex}
\mbox{}\hspace{0.11\tew}\footnotesize\phantom{\includegraphics[width=1.5cm]{logo_ctut}\hfill\raisebox{1ex}{Doãn Minh Đăng}}\hfill
\includegraphics[width=1.5cm]{logo_ctut}\hspace*{0.8cm}
\end{minipage}\\[-0.88\pah]

{\large #2}}
%-----------------------------------------------------------------------
%title CTUT and STW
\newcommand{\titstw}[1]{%\vspace*{-1ex}
{\hspace*{-0.0\tew}{\bf\large\color{black}{{#1}}}\\[-1.6ex] }\\[0.70\pah]
\fboxsep=0.3ex
\hspace*{-0.12\tew}
\colorbox{mag}{
\begin{minipage}{1.25\tew}
\tiny  \hfill \arabic{page}/\pageref{lastpage}\hspace*{1cm}
\end{minipage}}
\hspace*{-0.12\tew}
\begin{minipage}{1.25\tew}
\vspace*{-0.7ex}
\mbox{}\hspace{0.02\tew}
\raisebox{0.5ex}{\includegraphics[width=1.2cm]{logo_ctut}}
\hfill\footnotesize\raisebox{1.2ex}{Doãn Minh Đăng}\hfill
\raisebox{0.1ex}{\includegraphics[width=1.5cm]{logo_ctut}}\hspace*{0.8cm}
\end{minipage}\\[-0.88\pah]}
%-----------------------------------------------------------------------
%title CTUT logo small
\newcommand{\titls}[1]{%\vspace*{-1ex}
{\hspace*{-0.0\tew}{\bf\large\color{black}{{#1}}}\\[-1.6ex] }\\[0.70\pah]
\fboxsep=0.3ex
\hspace*{-0.12\tew}
\colorbox{mag}{
\begin{minipage}{1.25\tew}
\tiny  \hfill \arabic{page}/\pageref{lastpage}\hspace*{1cm}
\end{minipage}}
\hspace*{-0.12\tew}
\begin{minipage}{1.25\tew}
\vspace*{-0.7ex}
% \mbox{}\hspace{0.11\tew}\footnotesize\raisebox{1ex}{Carsten Scherer\ \ }\includegraphics[width=0.6cm]{logo_ctut}
% \hspace{26ex}\raisebox{1ex}{Siep Weiland\ \ }\includegraphics[width=0.6cm]{logo_ctut}\hspace*{0.8cm}
\mbox{}\hspace{0.11\tew}\footnotesize\raisebox{1ex}{Doãn Minh Đăng}\hfill
\includegraphics[width=0.5cm]{logo_ctut}\hspace*{0.8cm}
\end{minipage}\\[-0.88\pah]}
%-----------------------------------------------------------------------
%title underlined left
\newcommand{\titline}[1]{{
\hspace*{-0.0\tew}\hfill{\bf \color[rgb]{0,0,0}{#1}}\\[-1.6ex]}
\hspace*{-0.0\tew}\color[rgb]{0,0,0}{\rule{1.00\tew}{.12ex}}\\[-2ex]
\marginpar{{\tiny \vspace*{9cm}\hspace*{0.55cm}\arabic{page}/\pageref{lastpage}}}}

%-----------------------------------------------------------------------
%just title
\newcommand{\tit}[1]{
{\hspace*{-0.0\tew}{\bf\large \color{black}{{#1}}}\\[-1.6ex] }\\[0.79\pah]
\begin{minipage}{1.09\tew}
\tiny  \hfill \arabic{page}/\pageref{lastpage}
\end{minipage}\\[-0.9\pah]
}

%-----------------------------------------------------------------------
%title CTUT logo large
\newcommand{\titlogo}[2]{%\vspace*{-1ex}
{\hspace*{-0.0\tew}{\bf\Large\color{black}{{#1}}}\\[-1.6ex] }\\[0.70\pah]
\fboxsep=0.3ex
\hspace*{-0.12\tew}
\colorbox{mag}{
\begin{minipage}{1.25\tew}
\tiny  \hfill \arabic{page}/\pageref{lastpage}\hspace*{1cm}
\end{minipage}}
\hspace*{-0.12\tew}
\begin{minipage}{1.25\tew}
\vspace*{-0.7ex}
\mbox{}\hspace{0.11\tew}\footnotesize\raisebox{1ex}{Doãn Minh Đăng}\hfill
\includegraphics[width=0.6cm]{logo_ctut}\hspace*{0.8cm}
\end{minipage}\\[-0.88\pah]

{\large #2}}
%-----------------------------------------------------------------------
%title first page CTUT logo large
\newcommand{\titf}[2]{%\vspace*{-1ex}
{\hspace*{-0.0\tew}{\bf\Large\color{blue}{{#1}}}\\[-1.6ex] }\\[0.70\pah]
\fboxsep=0.3ex
\hspace*{-0.12\tew}
\colorbox{mag}{
\begin{minipage}{1.25\tew}
\tiny  \hfill \phantom{\arabic{page}/\pageref{lastpage}}\hspace*{1cm}
\end{minipage}}
\hspace*{-0.12\tew}
\begin{minipage}{1.25\tew}
\vspace*{-0.7ex}
\mbox{}\hspace{0.11\tew}\footnotesize\raisebox{1ex}{Doãn Minh Đăng}\hfill
\includegraphics[width=0.5cm]{logo_ctut}\hspace*{0.8cm}
\end{minipage}\\[-0.88\pah]

{\large #2}}

%-----------------------------------------------------------------------
% Boxes and minipages
%-----------------------------------------------------------------------
\newcommand{\sbo}[1]{
\shadowbox{
\vspace{1ex}
\begin{minipage}{0.97\tew}
\vspace*{1.5ex}
#1
\par\vspace{1ex}
\end{minipage}}}


\newcommand{\cpa}[3]{
{\fboxsep=1.2ex
\colorbox{#1}{
\begin{minipage}{#2}
#3
\end{minipage}}}
}

\newcommand{\cbo}[2]{
\fboxsep=1ex
\colorbox{#1}{
#2}}

\newcommand{\cbow}[3]{
\fboxsep=#2ex
\colorbox{#1}{
#3}}

\newcommand{\mini}[2]{ {\begin{minipage}[t]{#1}\sloppy #2 \end{minipage} } }

\newcommand{\twomini}[6]{
\setlength{\storee}{0.975\tew}
\addtolength{\storee}{-#1}\addtolength{\storee}{-#3}
\addtolength{\storee}{#2}
\mini{#1}{#5}\hspace{\storee}\raisebox{#4}{\mini{#3}{#6}} }
%\hspace{\storee\raisebox{#2}

\newcommand{\onemini}[3]{\hspace{#1}\mini{#2}{#3}}

%%%%%%%%%%%%%%%%%%%%%%%%%%%%%%%%%%%%%%%%%%%%%%%%%%%%%%%%%%%%%%%%%%%%%%%%%%%%%%%%%%%%%%%%%%%%%%%%%%%%%%%%%%%%%%%%%%%
%Colors
\newcommand{\red}[1]{{\color[rgb]{1,0,0}#1}}
\newcommand{\green}[1]{{\color[rgb]{0,0.7,0}#1}}
\newcommand{\blue}[1]{{\color[rgb]{0,0,1}#1}}
\newcommand{\dblue}[1]{{\color[rgb]{0,0,.4}#1}}
%\newcommand{\blue}[1]{{\color[rgb]{0,0.5,.8}#1}}
\newcommand{\black}[1]{{\color[rgb]{0,0,0}#1}}
%\newcommand{\yellow}[1]{{\color[rgb]{1,1,0}#1}}
\newcommand{\yellow}[1]{{\color[rgb]{1,0.95,.5}#1}}
%TUD RBG   0   174   239
\definecolor{lmag}{rgb}{0,0.82,1}
% \definecolor{mag}{rgb}{0,0.70,.88}
\newcommand{\tmag}[1]{{\color[rgb]{0,0.60,.78}#1}}
% \definecolor{mag}{rgb}{0,0.75,.93}
% \definecolor{mag}{rgb}{0,0.6,.8} %original
\definecolor{mag}{rgb}{0,0.7,.9}
\definecolor{lblue}{rgb}{0.6,0.8,1}
\definecolor{lgreen}{rgb}{0,0.92,0.7}
\definecolor{lred}{rgb}{1,0.75,0.75}
%\definecolor{lred}{rgb}{1,0.80,0.80}
\definecolor{lblue}{rgb}{0.6,0.8,1}
\definecolor{lgrey}{rgb}{0.85,0.85,0.85}

\definecolor{lye}{rgb}{1,0.9,0.6}
%\definecolor{lye}{rgb}{1,0.95,0.7}
%\definecolor{lgr}{rgb}{1,0.9,0.6}

\newcommand{\ypa}[1]{\cpa{lye}{0.95\tew}{#1}}
\newcommand{\bpa}[1]{\cpa{lblue}{0.95\tew}{#1}}
\newcommand{\rpa}[1]{\cpa{lred}{0.95\tew}{#1}}
%%%%%%%%%%%%%%%%%%%%%%%%%%%%%%%%%%%%%%%%%%%%%%%%%%%%%%%%%%%%%%%%%%%%%%%%%%%%%%%%%%%%%%%%%%%%%%%%%%%%%%%%%%%%%%%%%%%

\newcommand{\minimize}[2]{
\begin{array}{ll}
\mbox{minimize}    &\ \  #1 \\
\mbox{subject\ to} &\ \  #2
\end{array} }
\newcommand{\opt}{\text{opt}}

\newcommand{\X}{\red{X}}
\newcommand{\pa}{\blue{\delta}}
\newcommand{\ve}{\blue{v}}
\newcommand{\pas}{\blue{\bm{\delta}}}
\newcommand{\ves}{\bm{\blue{v}}}
\newcommand{\pad}{\blue{\bm{\dot{\delta}}}}
\newcommand{\pam}{\blue{\Delta(\pa)}}

\newcommand{\yc}{\blue{y}}
\newcommand{\cge}{\succcurlyeq}
\newcommand{\cle}{\preccurlyeq}
\newcommand{\cl}{\prec}
\newcommand{\cg}{\succ}
\newcommand{\rk}{\text{rank}}
\newcommand{\col}{\text{col}}

%\newcommand{\pa}{\blue{\delta}}
%\newcommand{\setpa}{\blue{\bm{\delta}}}

\renewcommand{\c}[1]{ {\cal #1} }
\newcommand{\bm}[1]{\text{\boldmath$ #1 $} }
\newcommand{\dis}{\displaystyle}
\newcommand{\ab}[1]{\langle #1 \rangle}
\newcommand{\ub}[2]{\underbrace{#1}_{#2}}

\newcommand{\diag}{ \operatornamewithlimits{diag} }
\newcommand{\re}{ \operatorname{Re} }
\newcommand{\he}{ \operatorname{He} }
\renewcommand{\t}[1]{ \tilde{\! #1\,} }
\newcommand{\tc}[1]{ \tilde{{\cal #1}} }
% \renewcommand{\h}[1]{ \hat{#1} } % 2014.08.06: remove this short-hand command, because it cause conflict with Vietnamese language
% If \h is used, some characters with accent like ể will cause error for amsmath: ``\hat allowed only in math mode''
\newcommand{\hc}[1]{ \hat{{\cal #1}} }
\newcommand{\esup}{ \operatornamewithlimits{ess\,sup} }
\newcommand{\tr}{ \operatorname{trace} }
\newcommand{\co}{ \mbox{co} }
\renewcommand{\i}{i}
\newcommand{\io}{i\omega}
\newcommand{\om}{\omega}
\newcommand{\la}{\lambda}
\newcommand{\ga}{\gamma}

\newcommand{\Cl}{ {\mathbb{C}^-} }
\newcommand{\Ci}{ {\mathbb{C}^0} }
\newcommand{\Cr}{ {\mathbb{C}^+} }
\newcommand{\C}{ {\mathbb{C}} }
\newcommand{\Cset}{ {\mathbb{C}} }

\newcommand{\Q}{ {\mathbb{Q}} }
\newcommand{\R}{ {\mathbb{R}} }
\newcommand{\Rset}{ {\mathbb{R}} }
\renewcommand{\H}{ {\mathbb{H}} }
\renewcommand{\S}{ {\mathbb{S}} }
\newcommand{\Nset}{ {\mathbb{N}} }
\newcommand{\Zset}{ {\mathbb{Z}} }
\newcommand{\alert}[1]{{\bf #1}}
%%%%%%%%%%%%%%%%%%%%%%%%%%%%%%%%%%%%%%%%%%%%%%%%%%%%%%%%%%%%%%%%%%%%%%%%%%%%%%%%%%%%%%%%%%%%%%%%%%%%%%%%%%%%%%%%%%%

\newcommand{\Alf}{{D}}
\newcommand{\Blf}{{C}}
\newcommand{\Clf}{{B}}
\newcommand{\Dlf}{{A}}


\newcommand{\Ac}{\green{A_c}}
\newcommand{\Bc}{\green{B_c}}
\newcommand{\Cc}{\green{C_c}}
\newcommand{\Dc}{\green{D_c}}

\newcommand{\Acl}{{\cal A}}
\newcommand{\Bcl}{{\cal B}}
\newcommand{\Ccl}{{\cal C}}
\newcommand{\Dcl}{{\cal D}}
\newcommand{\Xcl}{{ \green{\cal X} }}
\newcommand{\Ycl}{{\cal Y}}
\newcommand{\Zcl}{{\cal Z}}
\newcommand{\Pcl}{{\cal P}}
%closed-loop transformed
\newcommand{\A}{\bm{A}}
\newcommand{\B}{\bm{B}}
%\renewcommand{\C}{\bm{C}}
\newcommand{\D}{\bm{D}}
\renewcommand{\X}{\bm{X}}
\newcommand{\DX}{\bm{Y}}

\renewcommand{\v}{\red{v}}
\newcommand{\K}{\red{K}}
\renewcommand{\L}{\red{L}}
\newcommand{\M}{\red{M}}
\newcommand{\N}{\red{N}}
\renewcommand{\X}{\red{X}}
\newcommand{\Y}{\red{Y}}
\newcommand{\Z}{\red{Z}}
\newcommand{\F}{\blue{F}}

\newcommand{\T}{\c{T}}

\newcommand{\Ri}{\R\cup\{\infty\}}

%%%% Dang added %%%%
\newcommand{\xb}{\mathbf{x}}
\newcommand{\grad}{\nabla}
\newtheorem{Lem}{Bổ đề}
\newtheorem{Rem}{Lưu ý}
\newtheorem{Def}{Định nghĩa}
\newtheorem{Thm}{Định lý}
\newtheorem{Exa}{Ví dụ}



\documentclass[12pt]{article}
% \usepackage[utf8]{vietnam} % type Vietnamese
\usepackage[utf8]{inputenc}
\usepackage[vietnam]{babel}
\usepackage{amsmath,amsfonts,amssymb}
\usepackage{color,graphicx,fancybox} % pause.sty cung cap chuc nang pause giong nhu o documentclass beamer, tuy nhien dung o documentclass article nay thi khong co tac dung, nen bo di
\usepackage{sagetex} % run Sagemath within Latex
\usepackage{import} % 2014.08.05 so that using \import command instead of \input in the main files
% because \input or \include only search for the files with relative path to the MAIN tex files,
% while \import allows the figures inside the ``styles'' folder to be searched as well.
% see: http://sumitrangwala.wordpress.com/2009/03/09/latex-tips-and-tricks/

\graphicspath{{../styles/}} % khong co graphicspath thi file ../styles/slides.tex khong search duoc cac file logo, although da dung subimport
% Dung graphicspath la giai phap temporary, vi neu trong slides.tex ma co \input{otherfile.tex} thi no khong search duoc qua relative path
% Khi do phai kiem tra lam sao su dung duoc \subimport.

% \subimport*{../styles/}{slides} % dung thu ngay 2014.08.05 khong thay tac dung.

% 2014.08.05: This preamble block is now moved to each individual main file, because \import is now used instead of \input
% \documentclass[12pt]{article}
% \usepackage{amsmath,amssymb}
% \usepackage{color,graphicx,fancybox,pause}
% \usepackage[utf8]{vntex} % type Vietnamese
% \usepackage{sagetex} % run Sagemath within Latex

\renewcommand{\familydefault}{cmss}
\renewcommand{\floatpagefraction}{1}

%------------------------------------------
%  PDF
%------------------------------------------
\pdfoutput=1
\pdfcompresslevel=0
\pdfimageresolution=600
\pdfpkresolution=600
\pagestyle{empty}
\pdfpageheight=11.1cm
\pdfpagewidth=14.8cm
%------------------------------------------
% Text dimensions
%------------------------------------------
\textheight=12cm
\textwidth=12.2cm
\topmargin-3cm
\oddsidemargin-1.2cm
\parindent0cm
\parskip0ex
%\arraycolsep0.4ex
%------------------------------------------
%  Lengths
%------------------------------------------
\renewcommand{\baselinestretch}{1.2} % These two lines, although not called in .tex file, affect the output
\renewcommand{\arraystretch}{1}
\newcommand{\spa}[1]{\par\vspace*{#1ex}}
\newcommand{\skipthis}[1]{}

%Set standard display distances
%\rennewcommand{\arraystretch}{1.1}

\newcommand{\sdd}{
\setlength{\jot}{1ex}
\setlength{\abovedisplayskip}{1ex}
\setlength{\abovedisplayshortskip}{1ex}
\setlength{\belowdisplayskip}{1ex}
\setlength{\belowdisplayshortskip}{1ex}}

\newlength{\tew}\setlength{\tew}{\textwidth}
\newlength{\pah}\setlength{\pah}{\pdfpageheight}
\newlength{\storea}
\newlength{\storeb}
\newlength{\storec}
\newlength{\stored}
\newlength{\storee}
\newlength{\storef}

%Scale display distance box
\newcommand{\sdb}[2]{
\storea\abovedisplayskip
\storeb\abovedisplayshortskip
\storec\belowdisplayskip
\stored\belowdisplayshortskip
\setlength{\abovedisplayskip}{#1\abovedisplayskip}
\setlength{\abovedisplayshortskip}{#1\abovedisplayshortskip}
\setlength{\belowdisplayskip}{#1\belowdisplayskip}
\setlength{\belowdisplayshortskip}{#1\belowdisplayshortskip}

#2

\setlength{\abovedisplayskip}{\storea}
\setlength{\abovedisplayshortskip}{\storeb}
\setlength{\belowdisplayskip}{\storec}
\setlength{\belowdisplayshortskip}{\stored}
}

%Scale display distance box with different factors top and bottom
\newcommand{\sdbe}[3]{
\storea=\abovedisplayskip
\storeb=\abovedisplayshortskip
\storec=\belowdisplayskip
\stored=\belowdisplayshortskip
\setlength{\abovedisplayskip}{#1\abovedisplayskip}
\setlength{\abovedisplayshortskip}{#1\abovedisplayshortskip}
\setlength{\belowdisplayskip}{#2\belowdisplayskip}
\setlength{\belowdisplayshortskip}{#2\belowdisplayshortskip}

#3

\setlength{\abovedisplayskip}{\storea}
\setlength{\abovedisplayshortskip}{\storeb}
\setlength{\belowdisplayskip}{\storec}
\setlength{\belowdisplayshortskip}{\stored}
}

%Baselineskip box: scale baslineskip
\newcommand{\blb}[2]{
\storee=\baselineskip
\setlength{\baselineskip}{#1\baselineskip}
#2

\setlength{\baselineskip}{\storee}
}

%------------------------------------------
\newcommand{\enu}[1]{ \begin{enumerate} #1 \end{enumerate} }
\newcommand{\ite}[1]{\vspace*{0ex}\begin{list}{$\bullet$}{\leftmargin2.5ex}#1\end{list} } % itemize list without indentation
\newcommand{\itemi}[1]{\begin{itemize}#1\end{itemize} }

\renewcommand{\labelenumi}{\arabic{enumi}.}
\renewcommand{\labelenumii}{(\alph{enumii})}

\newcommand{\aln}[1]{\noindent\begin{align*}#1\end{align*}}
\newcommand{\mun}[1]{\noindent\begin{multline*}#1\end{multline*}}
\newcommand{\eqn}[1]{\spa{-3}$$#1$$\spa{0}}
\newcommand{\eqns}[1]{\spa{-2}$$#1$$\spa{0}}
\newcommand{\gan}[1]{\begin{gather*}#1\end{gather*}}
\newcommand{\ear}[1]{\begin{eqnarray}#1\end{eqnarray}}
\newcommand{\ean}[1]{\begin{eqnarray*}#1\end{eqnarray*}}

\newcommand{\arr}[2]{\begin{array}{#1}#2\end{array}}
\newcommand{\tab}[2]{\begin{tabular}{#1}#2\end{tabular}}
\newcommand{\mas}[2]{\left[\begin{array}{#1}#2\end{array}\right]}
\newcommand{\mat}[2]{\left(\begin{array}{#1}#2\end{array}\right)}
%\newcommand{\mat}[2]{\mas{#1}{#2}}
\newcommand{\matl}[4]{\left#1\begin{array}{#2}#3\end{array}\right#4}
\newcommand{\dmas}[2]{
\mas{ccc}{#1&&0\\[-1.5ex]&\ddots&\\[-1.5ex]0&&#2}}
\newcommand{\quo}[1]{\noindent\vspace*{-1ex}\begin{quote}#1\end{quote}}
\newcommand{\cen}[1]{\begin{center}#1\end{center}}
\newcommand{\epro}{ \mbox{\rule{2mm}{2mm}} }

%-----------------------------------------------------------------------
%title frontpage CTUT and STW
\newcommand{\titfstw}[2]{%\vspace*{-1ex}
{\hspace*{-0.0\tew}{\bf\Large\color{black}{{#1}}}\\[-1.6ex] }\\[0.70\pah]
\fboxsep=0.3ex
\hspace*{-0.12\tew}
\colorbox{mag}{
\begin{minipage}{1.25\tew}
\tiny  \hfill \phantom{\arabic{page}/\pageref{lastpage}}\hspace*{1cm}
\end{minipage}}
\hspace*{-0.12\tew}
\begin{minipage}{1.25\tew}
\vspace*{-0.7ex}
\mbox{}\hspace{0.11\tew}\footnotesize\phantom{\includegraphics[width=1.5cm]{logo_ctut}\hfill\raisebox{1ex}{Doãn Minh Đăng}}\hfill
\includegraphics[width=1.5cm]{logo_ctut}\hspace*{0.8cm}
\end{minipage}\\[-0.88\pah]

{\large #2}}
%-----------------------------------------------------------------------
%title CTUT and STW
\newcommand{\titstw}[1]{%\vspace*{-1ex}
{\hspace*{-0.0\tew}{\bf\large\color{black}{{#1}}}\\[-1.6ex] }\\[0.70\pah]
\fboxsep=0.3ex
\hspace*{-0.12\tew}
\colorbox{mag}{
\begin{minipage}{1.25\tew}
\tiny  \hfill \arabic{page}/\pageref{lastpage}\hspace*{1cm}
\end{minipage}}
\hspace*{-0.12\tew}
\begin{minipage}{1.25\tew}
\vspace*{-0.7ex}
\mbox{}\hspace{0.02\tew}
\raisebox{0.5ex}{\includegraphics[width=1.2cm]{logo_ctut}}
\hfill\footnotesize\raisebox{1.2ex}{Doãn Minh Đăng}\hfill
\raisebox{0.1ex}{\includegraphics[width=1.5cm]{logo_ctut}}\hspace*{0.8cm}
\end{minipage}\\[-0.88\pah]}
%-----------------------------------------------------------------------
%title CTUT logo small
\newcommand{\titls}[1]{%\vspace*{-1ex}
{\hspace*{-0.0\tew}{\bf\large\color{black}{{#1}}}\\[-1.6ex] }\\[0.70\pah]
\fboxsep=0.3ex
\hspace*{-0.12\tew}
\colorbox{mag}{
\begin{minipage}{1.25\tew}
\tiny  \hfill \arabic{page}/\pageref{lastpage}\hspace*{1cm}
\end{minipage}}
\hspace*{-0.12\tew}
\begin{minipage}{1.25\tew}
\vspace*{-0.7ex}
% \mbox{}\hspace{0.11\tew}\footnotesize\raisebox{1ex}{Carsten Scherer\ \ }\includegraphics[width=0.6cm]{logo_ctut}
% \hspace{26ex}\raisebox{1ex}{Siep Weiland\ \ }\includegraphics[width=0.6cm]{logo_ctut}\hspace*{0.8cm}
\mbox{}\hspace{0.11\tew}\footnotesize\raisebox{1ex}{Doãn Minh Đăng}\hfill
\includegraphics[width=0.5cm]{logo_ctut}\hspace*{0.8cm}
\end{minipage}\\[-0.88\pah]}
%-----------------------------------------------------------------------
%title underlined left
\newcommand{\titline}[1]{{
\hspace*{-0.0\tew}\hfill{\bf \color[rgb]{0,0,0}{#1}}\\[-1.6ex]}
\hspace*{-0.0\tew}\color[rgb]{0,0,0}{\rule{1.00\tew}{.12ex}}\\[-2ex]
\marginpar{{\tiny \vspace*{9cm}\hspace*{0.55cm}\arabic{page}/\pageref{lastpage}}}}

%-----------------------------------------------------------------------
%just title
\newcommand{\tit}[1]{
{\hspace*{-0.0\tew}{\bf\large \color{black}{{#1}}}\\[-1.6ex] }\\[0.79\pah]
\begin{minipage}{1.09\tew}
\tiny  \hfill \arabic{page}/\pageref{lastpage}
\end{minipage}\\[-0.9\pah]
}

%-----------------------------------------------------------------------
%title CTUT logo large
\newcommand{\titlogo}[2]{%\vspace*{-1ex}
{\hspace*{-0.0\tew}{\bf\Large\color{black}{{#1}}}\\[-1.6ex] }\\[0.70\pah]
\fboxsep=0.3ex
\hspace*{-0.12\tew}
\colorbox{mag}{
\begin{minipage}{1.25\tew}
\tiny  \hfill \arabic{page}/\pageref{lastpage}\hspace*{1cm}
\end{minipage}}
\hspace*{-0.12\tew}
\begin{minipage}{1.25\tew}
\vspace*{-0.7ex}
\mbox{}\hspace{0.11\tew}\footnotesize\raisebox{1ex}{Doãn Minh Đăng}\hfill
\includegraphics[width=0.6cm]{logo_ctut}\hspace*{0.8cm}
\end{minipage}\\[-0.88\pah]

{\large #2}}
%-----------------------------------------------------------------------
%title first page CTUT logo large
\newcommand{\titf}[2]{%\vspace*{-1ex}
{\hspace*{-0.0\tew}{\bf\Large\color{blue}{{#1}}}\\[-1.6ex] }\\[0.70\pah]
\fboxsep=0.3ex
\hspace*{-0.12\tew}
\colorbox{mag}{
\begin{minipage}{1.25\tew}
\tiny  \hfill \phantom{\arabic{page}/\pageref{lastpage}}\hspace*{1cm}
\end{minipage}}
\hspace*{-0.12\tew}
\begin{minipage}{1.25\tew}
\vspace*{-0.7ex}
\mbox{}\hspace{0.11\tew}\footnotesize\raisebox{1ex}{Doãn Minh Đăng}\hfill
\includegraphics[width=0.5cm]{logo_ctut}\hspace*{0.8cm}
\end{minipage}\\[-0.88\pah]

{\large #2}}

%-----------------------------------------------------------------------
% Boxes and minipages
%-----------------------------------------------------------------------
\newcommand{\sbo}[1]{
\shadowbox{
\vspace{1ex}
\begin{minipage}{0.97\tew}
\vspace*{1.5ex}
#1
\par\vspace{1ex}
\end{minipage}}}


\newcommand{\cpa}[3]{
{\fboxsep=1.2ex
\colorbox{#1}{
\begin{minipage}{#2}
#3
\end{minipage}}}
}

\newcommand{\cbo}[2]{
\fboxsep=1ex
\colorbox{#1}{
#2}}

\newcommand{\cbow}[3]{
\fboxsep=#2ex
\colorbox{#1}{
#3}}

\newcommand{\mini}[2]{ {\begin{minipage}[t]{#1}\sloppy #2 \end{minipage} } }

\newcommand{\twomini}[6]{
\setlength{\storee}{0.975\tew}
\addtolength{\storee}{-#1}\addtolength{\storee}{-#3}
\addtolength{\storee}{#2}
\mini{#1}{#5}\hspace{\storee}\raisebox{#4}{\mini{#3}{#6}} }
%\hspace{\storee\raisebox{#2}

\newcommand{\onemini}[3]{\hspace{#1}\mini{#2}{#3}}

%%%%%%%%%%%%%%%%%%%%%%%%%%%%%%%%%%%%%%%%%%%%%%%%%%%%%%%%%%%%%%%%%%%%%%%%%%%%%%%%%%%%%%%%%%%%%%%%%%%%%%%%%%%%%%%%%%%
%Colors
\newcommand{\red}[1]{{\color[rgb]{1,0,0}#1}}
\newcommand{\green}[1]{{\color[rgb]{0,0.7,0}#1}}
\newcommand{\blue}[1]{{\color[rgb]{0,0,1}#1}}
\newcommand{\dblue}[1]{{\color[rgb]{0,0,.4}#1}}
%\newcommand{\blue}[1]{{\color[rgb]{0,0.5,.8}#1}}
\newcommand{\black}[1]{{\color[rgb]{0,0,0}#1}}
%\newcommand{\yellow}[1]{{\color[rgb]{1,1,0}#1}}
\newcommand{\yellow}[1]{{\color[rgb]{1,0.95,.5}#1}}
%TUD RBG   0   174   239
\definecolor{lmag}{rgb}{0,0.82,1}
% \definecolor{mag}{rgb}{0,0.70,.88}
\newcommand{\tmag}[1]{{\color[rgb]{0,0.60,.78}#1}}
% \definecolor{mag}{rgb}{0,0.75,.93}
% \definecolor{mag}{rgb}{0,0.6,.8} %original
\definecolor{mag}{rgb}{0,0.7,.9}
\definecolor{lblue}{rgb}{0.6,0.8,1}
\definecolor{lgreen}{rgb}{0,0.92,0.7}
\definecolor{lred}{rgb}{1,0.75,0.75}
%\definecolor{lred}{rgb}{1,0.80,0.80}
\definecolor{lblue}{rgb}{0.6,0.8,1}
\definecolor{lgrey}{rgb}{0.85,0.85,0.85}

\definecolor{lye}{rgb}{1,0.9,0.6}
%\definecolor{lye}{rgb}{1,0.95,0.7}
%\definecolor{lgr}{rgb}{1,0.9,0.6}

\newcommand{\ypa}[1]{\cpa{lye}{0.95\tew}{#1}}
\newcommand{\bpa}[1]{\cpa{lblue}{0.95\tew}{#1}}
\newcommand{\rpa}[1]{\cpa{lred}{0.95\tew}{#1}}
%%%%%%%%%%%%%%%%%%%%%%%%%%%%%%%%%%%%%%%%%%%%%%%%%%%%%%%%%%%%%%%%%%%%%%%%%%%%%%%%%%%%%%%%%%%%%%%%%%%%%%%%%%%%%%%%%%%

\newcommand{\minimize}[2]{
\begin{array}{ll}
\mbox{minimize}    &\ \  #1 \\
\mbox{subject\ to} &\ \  #2
\end{array} }
\newcommand{\opt}{\text{opt}}

\newcommand{\X}{\red{X}}
\newcommand{\pa}{\blue{\delta}}
\newcommand{\ve}{\blue{v}}
\newcommand{\pas}{\blue{\bm{\delta}}}
\newcommand{\ves}{\bm{\blue{v}}}
\newcommand{\pad}{\blue{\bm{\dot{\delta}}}}
\newcommand{\pam}{\blue{\Delta(\pa)}}

\newcommand{\yc}{\blue{y}}
\newcommand{\cge}{\succcurlyeq}
\newcommand{\cle}{\preccurlyeq}
\newcommand{\cl}{\prec}
\newcommand{\cg}{\succ}
\newcommand{\rk}{\text{rank}}
\newcommand{\col}{\text{col}}

%\newcommand{\pa}{\blue{\delta}}
%\newcommand{\setpa}{\blue{\bm{\delta}}}

\renewcommand{\c}[1]{ {\cal #1} }
\newcommand{\bm}[1]{\text{\boldmath$ #1 $} }
\newcommand{\dis}{\displaystyle}
\newcommand{\ab}[1]{\langle #1 \rangle}
\newcommand{\ub}[2]{\underbrace{#1}_{#2}}

\newcommand{\diag}{ \operatornamewithlimits{diag} }
\newcommand{\re}{ \operatorname{Re} }
\newcommand{\he}{ \operatorname{He} }
\renewcommand{\t}[1]{ \tilde{\! #1\,} }
\newcommand{\tc}[1]{ \tilde{{\cal #1}} }
% \renewcommand{\h}[1]{ \hat{#1} } % 2014.08.06: remove this short-hand command, because it cause conflict with Vietnamese language
% If \h is used, some characters with accent like ể will cause error for amsmath: ``\hat allowed only in math mode''
\newcommand{\hc}[1]{ \hat{{\cal #1}} }
\newcommand{\esup}{ \operatornamewithlimits{ess\,sup} }
\newcommand{\tr}{ \operatorname{trace} }
\newcommand{\co}{ \mbox{co} }
\renewcommand{\i}{i}
\newcommand{\io}{i\omega}
\newcommand{\om}{\omega}
\newcommand{\la}{\lambda}
\newcommand{\ga}{\gamma}

\newcommand{\Cl}{ {\mathbb{C}^-} }
\newcommand{\Ci}{ {\mathbb{C}^0} }
\newcommand{\Cr}{ {\mathbb{C}^+} }
\newcommand{\C}{ {\mathbb{C}} }
\newcommand{\Cset}{ {\mathbb{C}} }

\newcommand{\Q}{ {\mathbb{Q}} }
\newcommand{\R}{ {\mathbb{R}} }
\newcommand{\Rset}{ {\mathbb{R}} }
\renewcommand{\H}{ {\mathbb{H}} }
\renewcommand{\S}{ {\mathbb{S}} }
\newcommand{\Nset}{ {\mathbb{N}} }
\newcommand{\Zset}{ {\mathbb{Z}} }
\newcommand{\alert}[1]{{\bf #1}}
%%%%%%%%%%%%%%%%%%%%%%%%%%%%%%%%%%%%%%%%%%%%%%%%%%%%%%%%%%%%%%%%%%%%%%%%%%%%%%%%%%%%%%%%%%%%%%%%%%%%%%%%%%%%%%%%%%%

\newcommand{\Alf}{{D}}
\newcommand{\Blf}{{C}}
\newcommand{\Clf}{{B}}
\newcommand{\Dlf}{{A}}


\newcommand{\Ac}{\green{A_c}}
\newcommand{\Bc}{\green{B_c}}
\newcommand{\Cc}{\green{C_c}}
\newcommand{\Dc}{\green{D_c}}

\newcommand{\Acl}{{\cal A}}
\newcommand{\Bcl}{{\cal B}}
\newcommand{\Ccl}{{\cal C}}
\newcommand{\Dcl}{{\cal D}}
\newcommand{\Xcl}{{ \green{\cal X} }}
\newcommand{\Ycl}{{\cal Y}}
\newcommand{\Zcl}{{\cal Z}}
\newcommand{\Pcl}{{\cal P}}
%closed-loop transformed
\newcommand{\A}{\bm{A}}
\newcommand{\B}{\bm{B}}
%\renewcommand{\C}{\bm{C}}
\newcommand{\D}{\bm{D}}
\renewcommand{\X}{\bm{X}}
\newcommand{\DX}{\bm{Y}}

\renewcommand{\v}{\red{v}}
\newcommand{\K}{\red{K}}
\renewcommand{\L}{\red{L}}
\newcommand{\M}{\red{M}}
\newcommand{\N}{\red{N}}
\renewcommand{\X}{\red{X}}
\newcommand{\Y}{\red{Y}}
\newcommand{\Z}{\red{Z}}
\newcommand{\F}{\blue{F}}

\newcommand{\T}{\c{T}}

\newcommand{\Ri}{\R\cup\{\infty\}}

%%%% Dang added %%%%
\newcommand{\xb}{\mathbf{x}}
\newcommand{\grad}{\nabla}
\newtheorem{Lem}{Bổ đề}
\newtheorem{Rem}{Lưu ý}
\newtheorem{Def}{Định nghĩa}
\newtheorem{Thm}{Định lý}
\newtheorem{Exa}{Ví dụ}



%% end of slides.tex
% % \renewcommand{\pause}{}
% \newcommand{\Hi}{H_\infty}

\begin{document}
% \def\firstofone#1{#1}
% \let\uppercase\firstofone
% \let\MakeUppercase\firstofone

\sdd

% \newpage\titlogo{Tiêu đề có logo lớn}

% \newpage\titf{Tiêu đề có logo khá lớn, dùng làm trang đầu}

% \newpage\titls{Tiêu đề có logo nhỏ}

% spa{1} Tạo một khoảng cách giữa 2 dòng cho dễ nhìn

\newpage\titf{Chương 4: Nội suy và xấp xỉ hàm}

\enu{
\item Khái niệm đa thức nội suy
\item Đa thức nội suy Lagrange
\item Đa thức nội suy Newton
\item Bài toán xấp xỉ hàm thực nghiệm và phương pháp bình phương bé nhất
% \item Bài tập
% \item Lập trình giải thuật
}

\newpage\titls{Khái niệm đa thức nội suy}

Giả sử ta có một mối quan hệ giữa hai đại lượng, viết ở dạng hàm số: $y=f(x)$, tuy nhiên ta \textbf{chưa biết} công thức của hàm số $f(x)$ mà chỉ có các cặp giá trị $(x_k,y_k)$ ở một số hữu hạn $n+1$ điểm ($k=\overline{0,n}$). Các giá trị $x_k$ được gọi là \textbf{điểm nút} và xếp thứ tự tăng dần: $x_0 < x_1 < \cdots < x_n$.
\spa{1}
\blue{Vấn đề được quan tâm: xây dựng đa thức $P(x)$ thỏa mãn $P(x_k)=y_k, k=\overline{0,n}$. Ta gọi $P(x)$ là đa thức nội suy của hàm số $f(x)$.}
\spa{1}
\textbf{Lưu ý}: có nhiều đa thức $P(x)$ cùng thỏa điều kiện đi qua bộ các điểm $(x_k,y_k), k=\overline{0,n}$, tuy nhiên ta chứng minh được chỉ tồn tại duy nhất một đa thức nội suy $P(x)$ có bậc nhỏ hơn hoặc bằng $n$.

\newpage\titls{Khái niệm đa thức nội suy}

\begin{sagesilent}
 f=-x^2*sin(x) + 2
 n=5
 xk=[2*i for i in range(n)]
 fxk=range(n)
 yk=[]
 for i in range(n):
   fxk[i]=N(f(x=xk[i]),20)
   # Neu khong chuyen thanh so thuc bang N(.,20) thi fxk[i] la bieu thuc symbolic dai, dua vao tableVD1 se bao loi
   yk.append((xk[i],fxk[i]))
 # Dung lenh sau tao ra list nhanh hon
 # yk2=[(xk[i],fxk[i]) for i in range(n)]
 R = PolynomialRing(QQ, 'x')
 Lagrange_1 = R.lagrange_polynomial(yk)
 plotL1=plot(Lagrange_1, (-1),(9), color='red', legend_label='$P(x)$')

 plotfxk=list_plot(yk, size=50, legend_label='($x_k$,$y_k$)')
 tableVD1=r"\begin{tabular}{l|c|l}"
 tableVD1+=r"$k$ & $x_k$ & $y_k$ \\ \hline"
 for i in range(n):
   tableVD1+=latex(i) + r"&" + latex(round(xk[i],2)) + r"&" + latex(fxk[i]) + r"\\"
 tableVD1+=r"\end{tabular}"
\end{sagesilent}

Ví dụ:

%    \cen{\sagestr{tableVD1}}
   
 \begin{tabular}{lc}
   \sagestr{tableVD1}
  & 
    \begin{tabular}{c}
    \sageplot[scale=0.4]{plotL1+plotfxk}
    \end{tabular}
 \end{tabular}


\newpage\titls{Đa thức nội suy Lagrange}

Ta tìm đa thức có bậc nhỏ hơn hay bằng $n$ của hàm số $f(x)$ trên $[x_0, x_n]$, gọi là \textbf{đa thức nội suy Lagrange}.

\blue{Ý tưởng: tạo ra các \textbf{đa thức cơ sở} bằng với $1$ tại điểm $x_k$ và bằng $0$ tại các điểm $x_j, j \neq k$, để dễ dàng đạt được giá trị $y_k$ tại mỗi điểm nút.}

Ta xây dựng các đa thức $p_n^{(k)}(x)$ có bậc bằng $n$ và thỏa điều kiện:
\begin{align}\label{eq_Lagrange_pk_aim}
 p_n^{(k)} \big( x_j \big) = \left\{ \begin{array}{lc}
                                      1 & \qquad j=k \\
                                      0 & \qquad j \neq k
                                     \end{array}
                                     \right.
\end{align}

Từ các đa thức $p_n^{(k)}(x)$, ta lập đa thức nội suy Lagrange:
\red{\begin{align}\label{eq_Lagrange_L}
\mathbf{\mathcal{L}_n(x) = \sum_{k=0}^n p_n^{(k)}(x).y_k}
\end{align}
}
Do \eqref{eq_Lagrange_pk_aim}, đa thức $\mathcal{L}_n(x)$ theo công thức \eqref{eq_Lagrange_L} thỏa mãn yêu cầu $\mathcal{L}_n(x_k)=y_k, \forall k=\overline{0,n}$, và có bậc nhỏ hơn hoặc bằng $n$.

\newpage\titls{Đa thức nội suy Lagrange}

\blue{Vấn đề: lập các đa thức $p_n^{(k)}$ như thế nào để thỏa điều kiện \eqref{eq_Lagrange_pk_aim}?}

Đa thức $p_n^{(k)}(x)$ bậc $n$ và có $n$ nghiệm $x_0, \cdots, x_{k-1}, x_{k+1}, x_n$, nên ta có thể viết dưới dạng:
\begin{align}\label{eq_Lagrange_pk}
 p_n^{(k)}(x) = C_k(x-x_0) \cdots (x-x_{k-1})(x-x_{k+1}) \cdots (x-x_{n})
\end{align}
với hằng số $C_k$.

Để $p_n^{(k)}(x_k)=1$ thì thay $x=x_k$ vào công thức \eqref{eq_Lagrange_pk}, ta tìm được $C_k$:
\begin{align}
 C_k=\frac{1}{(x_k-x_0) \cdots (x_k-x_{k-1})(x_k-x_{k+1}) \cdots (x_k-x_{n})}
\end{align}

Vậy ta tìm được công thức của đa thức $p_n^{(k)}(x)$:
\red{\begin{align}
\mathbf{p_n^{(k)}(x)=\frac{(x-x_0) \cdots (x-x_{k-1})(x-x_{k+1}) \cdots (x-x_{n})}{(x_k-x_0) \cdots (x_k-x_{k-1})(x_k-x_{k+1}) \cdots (x_k-x_{n})}}
\end{align}
}

\newpage\titls{Đa thức nội suy Lagrange}

\begin{sagesilent}
 n=3
 xk2=[0,1,3]
 fxk2=[1,-1,2]
 #yk2=[(xk2[i],fxk2[i]) for i in range(n)]
 yk2=zip(xk2,fxk2) # zip: couple pairs of values of the two arrays
 R = PolynomialRing(QQ, 'x')
 Lagrange_2 = R.lagrange_polynomial(yk2)
 plotL2=plot(Lagrange_2, (-1),(4), color='red', legend_label='$\mathcal{L}_2(x)$')
 # Ref: http://www.sagemath.org/doc/reference/polynomial_rings/sage/rings/polynomial/polynomial_ring.html
 #Lagrange_2=7/6*x^2-19/6*x+1
 #plotL2=plot(Lagrange_2, (-1,4), color='red', legend_label='$\mathcal{L}_2(x)$')
 #plotfxk2=list_plot(yk2, size=50, legend_label='($x_k$,$y_k$)')
 plotfxk2=list_plot(yk2, size=50)
 tableVD2=r"\begin{tabular}{l|c|l}"
 tableVD2+=r"$k$ & $x_k$ & $y_k$ \\ \hline"
 for i in range(n):
   tableVD2+=latex(i) + r"&" + latex(round(xk2[i],2)) + r"&" + latex(fxk2[i]) + r"\\"
 tableVD2+=r"\end{tabular}"
\end{sagesilent}

Ví dụ: tính đa thức nội suy Lagrange với các giá trị cho trong bảng số:

 \begin{tabular}{lc}
   \sagestr{tableVD2}
  & 
    \begin{tabular}{c}
    \sageplot[scale=0.4]{plotfxk2}
    \end{tabular}
 \end{tabular}

\newpage\titls{Đa thức nội suy Lagrange}

Kết quả: $\mathcal{L}_2(x)=\frac{7}{6}x^2 - \frac{19}{6}x + 1$

 \begin{tabular}{lc}
   \sagestr{tableVD2}
  & 
    \begin{tabular}{c}
    \sageplot[scale=0.4]{plotfxk2+plotL2}
    \end{tabular}
 \end{tabular}

\newpage\titls{Đa thức nội suy Lagrange}
 
Cách dễ nhớ để tính các hệ số của đa thức nội suy Lagrange: lập bảng.

\begin{tabular}{c|cccc|c}
 $x$ & $x_0$ & $x_1$ & $\cdots$ & $x_n$ & \\ \hline
 $x_0$ & $\mathbf{x-x_0}$ & $x_0-x_1$ & $\cdots$ & $x_0-x_n$ & $D_0$ \\
 $x_1$ & $x_1-x_0$ & $\mathbf{x-x_1}$ & $\cdots$ & $x_1-x_n$ & $D_1$ \\
 $\cdots$ & $\cdots$ & $\cdots$ & $\cdots$ & $\cdots$ & $\cdots$ \\
 $x_n$ & $x_n-x_0$ & $x_n-x_1$ & $\cdots$ & $\mathbf{x-x_n}$ & $D_n$ \\ \hline
 & & & & & $\mathbf{w(x)}$
\end{tabular}

Lấy tích theo hàng ngang thứ $k$ là đa thức $D_k$, còn lấy tích theo đường chéo chính thì ta được $w(x)$. Khi đó đa thức nội suy Lagrange là:
\red{\begin{align}\label{eq_Lagrange_L_type2}
\mathbf{\mathcal{L}_n(x) = w(x) \sum_{k=0}^n \frac{y_k}{D_k}}
\end{align}
}

Ngoài ra, với mỗi điểm $x \in [x_0,x_n]$ bất kỳ, ta thay giá trị đó vào bảng và dùng công thức \eqref{eq_Lagrange_L_type2} thì tìm ra được giá trị nội suy $\mathcal{L}_n(x)$.

\newpage\titls{Đa thức nội suy Lagrange}

Ví dụ: tìm giá trị nội suy tại điểm $x=2$ từ bảng số:

   \cen{\sagestr{tableVD2}}
% Sử dụng lại ví dụ số 2
\begin{sagesilent}
xp=2 
\end{sagesilent}

\enu{
\item Dùng công thức vừa tính được ở ví dụ trước: $\mathcal{L}_2(x)=\frac{7}{6}x^2 - \frac{19}{6}x + 1$ 
\item Lập bảng để vừa xây dựng công thức, vừa tính giá trị nội suy
}

\newpage\titls{Đa thức nội suy Lagrange}

\cen{\begin{tabular}{lr}
\sagestr{tableVD2}
  &
 \begin{tabular}{c|rrr|r}
 $x=\sage{xp}$ & $\sage{xk2[0]}$ & $\sage{xk2[1]}$ & $\sage{xk2[2]}$ & \\ \hline
 $\sage{xk2[0]}$ & $\sage{xp-xk2[0]}$ & $\sage{xk2[0]-xk2[1]}$ & $\sage{xk2[0]-xk2[2]}$ & $\sage{(xp-xk2[0])*(xk2[0]-xk2[1])*(xk2[0]-xk2[2])}$ \\
 $\sage{xk2[1]}$ & $\sage{xk2[1]-xk2[0]}$ & $\sage{xp-xk2[1]}$ & $\sage{xk2[1]-xk2[2]}$ & $\sage{(xk2[1]-xk2[0])*(xp-xk2[1])*(xk2[1]-xk2[2])}$ \\
 $\sage{xk2[2]}$ & $\sage{xk2[2]-xk2[0]}$ & $\sage{xk2[2]-xk2[1]}$ & $\sage{xp-xk2[2]}$ & $\sage{(xk2[2]-xk2[0])*(xk2[2]-xk2[1])*(xp-xk2[2])}$ \\ \hline
 & & & & $\sage{(xp-xk2[0])*(xp-xk2[1])*(xp-xk2[2])}$
 \end{tabular}
\end{tabular}
}

Tính giá trị nội suy tại $x=2$:
\begin{align}
 \mathcal{L}_2(2)=\sage{(xp-xk2[0])*(xp-xk2[1])*(xp-xk2[2])}\Bigg( \frac{\sage{fxk2[0]}}{\sage{(xp-xk2[0])*(xk2[0]-xk2[1])*(xk2[0]-xk2[2])}} + \frac{\sage{fxk2[1]}}{\sage{(xk2[1]-xk2[0])*(xp-xk2[1])*(xk2[1]-xk2[2])}} + \frac{\sage{fxk2[2]}}{\sage{(xk2[2]-xk2[0])*(xk2[2]-xk2[1])*(xp-xk2[2])}} \Bigg) = \sage{(xp-xk2[0])*(xp-xk2[1])*(xp-xk2[2])*(fxk2[0]/((xp-xk2[0])*(xk2[0]-xk2[1])*(xk2[0]-xk2[2])) + fxk2[1]/((xk2[1]-xk2[0])*(xp-xk2[1])*(xk2[1]-xk2[2])) + fxk2[2]/((xk2[2]-xk2[0])*(xk2[2]-xk2[1])*(xp-xk2[2])))}
\end{align}

\newpage\titls{Đa thức nội suy Lagrange}

\begin{sagesilent}
 n=4
 xk3=[0,1,3,4]
 fxk3=[1,1,2,-1]
# yk3=[(xk3[i],fxk3[i]) for i in range(n)]
 yk3=zip(xk3,fxk3) # zip: couple pairs of values of the two arrays
 R = PolynomialRing(QQ, 'x')
 Lagrange_3 = R.lagrange_polynomial(yk3)
 plotL3=plot(Lagrange_3, (-1),(5), color='red', legend_label='$\mathcal{L}_3(x)$')
 plotfxk3=list_plot(yk3, size=50, legend_label='($x_k$,$y_k$)')
 tableVD3=r"\begin{tabular}{l|c|l}"
 tableVD3+=r"$k$ & $x_k$ & $y_k$ \\ \hline"
 for i in range(n):
   tableVD3+=latex(i) + r"&" + latex(round(xk3[i],2)) + r"&" + latex(fxk3[i]) + r"\\"
 tableVD3+=r"\end{tabular}"
\end{sagesilent}

Ví dụ: dùng nội suy Lagrange tính giá trị tại $x=2$ với các giá trị cho trong bảng số:

 \begin{tabular}{lc}
   \sagestr{tableVD3}
  & 
    \begin{tabular}{c}
    \sageplot[scale=0.4]{plotfxk3+plotL3}
    \end{tabular}
 \end{tabular}

 \newpage\titls{Đa thức nội suy Lagrange}

 Lập bảng:
 
 \begin{tabular}{c|rrrr|r}
 $x=2$ & 0 & 1 & 3 & 4 & \\ \hline
 0 & 2 & -1 & -3 & -4 & -24 \\
 1 & 1 & 1 & -2 & -3 & 6 \\
 3 & 3 & 2 & -1 & -1 & 6 \\
 4 & 4 & 3 & 1 & -2 & -24 \\ \hline
 & & & & & 4
\end{tabular}

 Tính giá trị nội suy:
\begin{align*}
 \mathcal{L}_3(2)=4\Bigg( \frac{1}{24} + \frac{1}{6} + \frac{2}{6} + \frac{-1}{24}\Bigg) = 2
\end{align*}

Công thức của đa thức nội suy: $\mathcal{L}_3(x)=\sage{Lagrange_3}$

\newpage\titls{Đa thức nội suy Newton}

\blue{Câu hỏi: có cách nào khác để lập bảng và tính toán giá trị của đa thức nội suy tại một điểm $x^*$ không?

Trả lời: ta có thể viết đa thức nội suy dưới dạng \textbf{đa thức nội suy Newton}, và lập bảng để tính.}

Đa thức nội suy Newton được định nghĩa dựa trên \textbf{sai phân hữu hạn}: $\Delta_{y_k}=y_{k+1}-y_k, \Delta_{x_k}=x_{k+1}-x_k$.

Ta giới thiệu khái niệm \textbf{tỉ sai phân cấp 1} của hàm $f$ trên doạn $[x_k,x_{k+1}]$:
\begin{equation}
 f[x_k,x_{k+1}]=\frac{\Delta_{y_k}}{\Delta_{x_k}}=\frac{y_{k+1}-y_k}{x_{k+1}-x_k}
\end{equation}
và \textbf{tỉ sai phân cấp} $\mathbf{p}$ (áp dụng cho $p=2,3,\cdots$):
\begin{equation}
 f[x_k,\cdots,x_{k+p}]=\frac{f[x_{k+1},\cdots,x_{k+p}]-f[x_{k},\cdots,x_{k+p-1}]}{x_{k+p}-x_k}
\end{equation}

\newpage\titls{Đa thức nội suy Newton}

Công thức của đa thức nội suy Newton:
\begin{align}\label{eq_Newton}
 f(x)&=y_0 + f[x_0,x_1](x-x_0) + f[x_0,x_1,x_2](x-x_0)(x-x_1) + \cdots + \nonumber \\
 &+f[x_0,\cdots,x_n](x-x_0)\cdots (x-x_{n-1}) + \nonumber \\
 &+f[x,x_0,\cdots,x_n](x-x_0)\cdots (x-x_{n-1})(x-x_{n})
\end{align}

Ta đặt $\mathcal{R}_n(x)$ là số hạng cuối của \eqref{eq_Newton}, và phần trước là $\mathcal{N}_n^f(x)$:
\red{\begin{align}
 \mathcal{N}_n^f(x)&=y_0 + f[x_0,x_1](x-x_0) + f[x_0,x_1,x_2](x-x_0)(x-x_1) + \cdots + \nonumber \\
 &+f[x_0,\cdots,x_n](x-x_0)\cdots (x-x_{n-1}) \label{eq_Newton_forward}\\
 \mathcal{R}_n(x)&=f[x,x_0,\cdots,x_n](x-x_0)\cdots (x-x_{n-1})(x-x_{n}) \label{eq_Newton_res}
\end{align}
}
Khi đó ta viết gọn lại:
\red{\begin{align}\label{eq_Newton2}
 f(x)=\mathcal{N}_n^f(x)+\mathcal{R}_n(x) \approx \mathcal{N}_n^f(x)
\end{align}
}

\newpage\titls{Đa thức nội suy Newton}

Để tính toán giá trị hàm số $f(x)$ theo công thức \eqref{eq_Newton2}, tức  $f(x)\approx \mathcal{N}_n^f(x)$ với sai số là $\mathcal{R}_n(x)$, ta lập \textbf{bảng tỉ sai phân}.

\begin{sagesilent}
 xp=1.5
 n=6
 xk4=[1,2,3,4,5,6]
 fxk4=[-3,0,15,48,105,192]
 #[(1, -3), (2, 0), (3, 15), (4, 48), (5, 105), (6, 192)]
 yk4=zip(xk4,fxk4) # zip: couple pairs of values of the two arrays
 R = PolynomialRing(RR, "x")
 bangNewton = R.divided_difference(yk4, full_table=True)
 Lagrange_4 = R.lagrange_polynomial(yk4)
 plotL4=plot(Lagrange_4, (-1),(7), color='red', legend_label='$\mathcal{L}_5(x)$')
 plotfxk4=list_plot(yk4, size=50, legend_label='($x_k$,$y_k$)')
 tableVD4=r"\begin{tabular}{l|c|l}"
 tableVD4+=r"$k$ & $x_k$ & $y_k$ \\ \hline"
 for i in range(n):
   tableVD4+=latex(i) + r"&" + latex(round(xk4[i],2)) + r"&" + latex(fxk4[i]) + r"\\"
 tableVD4+=r"\end{tabular}"
\end{sagesilent}

\begin{sagesilent}
# Now change the full table from Sagemath to my form of representing table
 bang2=list(list(i for i in range(j,n)) for j in range(n))
 for j in range(n):            
    for i in range(j,n):
        bang2[j][i-j]=bangNewton[i][j]
# Now create the Latex table
 tableNewton=r'\begin{tabular}{l|c|' # first column: k, second column: x_k
 for i in range(n): tableNewton+='c' # n columns for finit difference order 0 to n
 tableNewton+=r'}'
 tableNewton+=r'$k$ & $x_k$ & $y_k$'
# for i in range(n-1): tableNewton+=r' & f[$x_0,\cdots,x_'+str(i+1)+r'$]'
 for i in range(1,n): tableNewton+=r' & $f_\Delta^'+str(i)+r'$'
 tableNewton+=r'\\ \hline'
 for i in range(n):
   tableNewton+=latex(i) + r' & ' + latex(round(xk4[i],2))
   for j in range(i+1):
     tableNewton+= r' & ' + latex(round(bangNewton[i][j],2))
   for j in range(i+1,n): tableNewton+= r' & '
   tableNewton+= r'\\'
 tableNewton+=r'\end{tabular}'
\end{sagesilent}

Ví dụ: Tính giá trị gần đúng của hàm số tại $x=\sage{round(xp,2)}$, với dữ liệu:

 \begin{tabular}{lc}
   \sagestr{tableVD4}
  & 
    \begin{tabular}{c}
    \sageplot[scale=0.4]{plotfxk4+plotL4}
    \end{tabular}
 \end{tabular}


\newpage\titls{Đa thức nội suy Newton}

Kết quả tính bảng tỉ sai phân: (ký hiệu $f_\Delta^k=f[x_0,\cdots,x_k]$)

\sagestr{tableNewton}

% \begin{equation*}
%  \sage{bangNewton}
% \end{equation*}

\begin{sagesilent}
 mathNewton = r' = '+latex(bangNewton[0][0])
 sohangNewton = [bangNewton[i][i] for i in range(n)]
 for i in range(1,n):
   mathNewton+= r' + '+latex(bangNewton[i][i])
   for j in range(i): 
     mathNewton+=r'('+latex(round(xp-xk4[j],2))+r')'
     sohangNewton[i] *= xp-xk4[j] 
 ketquaNewton = sum(sohangNewton[i] for i in range(n))
\end{sagesilent}

\begin{align*}
 f(\sage{round(xp,2)}) &\approx \mathcal{N}_{\sage{n}}^f(\sage{round(xp,2)}) \\
 &\sagestr{mathNewton} \\
 &=\sage{round(ketquaNewton,3)}
\end{align*}


\newpage\titls{Đa thức nội suy Newton}

Vài lưu ý:
\itemi{
\item Biểu thức $\mathcal{R}_n(x)$ (bậc $n+1$) là sai số của đa thức nội suy Newton.
\item Công thức \eqref{eq_Newton_forward} gọi là \textbf{công thức Newton tiến (forward)}, ta còn có thể dùng \textbf{công thức Newton lùi (backward)}:
\begin{align} \label{eq_Newton_backward}
 \mathcal{N}_n^b(x)&=y_n + f[x_{n-1},x_{n}](x-x_n) + \nonumber \\
 &+f[x_{n-1},x_{n-1},x_n](x-x_{n-1})(x-x_{n-2}) + \cdots + \nonumber \\
 &+f[x_0,\cdots,x_n](x-x_1)\cdots (x-x_{n})
 \end{align}
\item Do đa thức nội suy bậc $n$ trở xuống là duy nhất (khi đi qua bộ $n+1$ điểm), nên khi xây dựng bằng phương pháp Lagrange và Newton thì sẽ ra kết quả giống nhau: \red{$\mathcal{L}_n(x)=\mathcal{N}_n^f(x)=\mathcal{N}_n^b(x)$}.
% \begin{equation}
%   \mathcal{L}_n(x)=\mathcal{N}_n^f(x)=\mathcal{N}_n^b(x)
%  \end{equation}
 }

\newpage\titls{Đa thức nội suy spline bậc ba}

\blue{Vấn đề: Nếu có nhiều điểm, phải lập đa thức nội suy bậc cao, phải tính toán nhiều. Có dạng đa thức nội suy nào bậc thấp, linh hoạt hơn?}
 
 \textbf{Định nghĩa spline:}
 
 Một spline bậc ba nội suy của một hàm $f(x)$ trên đoạn $[a,b]$ là hàm $g(x)$ thỏa các điều kiện:
 \itemi{
 \item $g(x)$ có \textbf{đạo hàm đến cấp hai liên lục} trên đoạn $[a,b]$.
 \item Trên mỗi đoạn $[x_k,x_{k+1}], k=0,\cdots, n-1$, ta có $g(x)=g_k(x)$ là một đa thức bậc ba.
 \item $g(x_k)=f(x_k)=y_k, \forall k$.
}
\blue{Tóm lại: spline là đa thức bậc ba trên mỗi khoảng $[x_k,x_{k+1}], k=0,\cdots, n-1$, và có đặc điểm là liên tục, khả vi tại các điểm nối.}

\newpage\titls{Bài toán xấp xỉ hàm thực nghiệm}

\blue{Một vấn đề khác: Giả sử có một tập hợp điểm là từ một hàm đơn giản (ví dụ: đường thẳng), nhưng do có các \textbf{nhiễu} nên có thể nằm lệch. Làm sao tìm một hàm số theo dạng đã biết và \textbf{gần} tập hợp điểm đó nhất? (không nhất thiết phải đi qua tất cả các điểm)

Cách giải quyết: \textbf{Phương pháp bình phương bé nhất} (hoặc: PP bình phương cực tiểu, PP bình phương tối thiểu).}

Cho tập hợp điểm $(x_k,y_k), k=1,\cdots,n$, ta tìm một hàm $f(x)$ sao cho tổng sau là nhỏ nhất:
\begin{align}
g(f)= \sum_{k=1}^n [f(x_k)-y_k]^2
\end{align}

\newpage\titls{Bài toán xấp xỉ hàm thực nghiệm}

Trường hợp thường gặp: hàm $f(x)$ là hàm số tuyến tính (đa thức bậc một), tức là $f(x)=A+Bx$. Bài toán phát biểu là:

\begin{align}
 \min_{A,B} g(A,B) = \min_{A,B} \sum_{k=1}^n (A+Bx_k-y_k)^2
\end{align}

\blue{Cách giải: để tìm giá trị $A^*,B^*$ mà $g(A^*,B^*)=\min_{A,B} g(A,B)$, ta tìm đạo hàm của $g$ theo các biến $A,B$, và giải để đạo hàm bằng không:}

\begin{align}
\left\lbrace\begin{array}{rcccl}
 \frac{\partial g}{\partial A} &=& 2 \sum_{k=1}^n (A+Bx_k-y_k) &=& 0 \\
 \frac{\partial g}{\partial B} &=& 2 \sum_{k=1}^n (A+Bx_k-y_k)x_k &=& 0
\end{array}
\right.
\end{align}

\newpage\titls{Bài toán xấp xỉ hàm thực nghiệm}

\blue{
Ta viết lại thành hệ phương trình tuyến tính với hai ẩn là $A$ và $B$:
\begin{align}\label{eq_linear_leastsquare}
\left\lbrace\begin{array}{rcccl}
 n \mathbf{A} &+& \left(\sum_{k=1}^n x_k \right) \mathbf{B} &=& \sum_{k=1}^n y_k \\
 \left(\sum_{k=1}^n x_k\right) \mathbf{A} &+& \left(\sum_{k=1}^n x_k^2 \right) \mathbf{B} &=& \sum_{k=1}^n x_k y_k
\end{array}
\right.
\end{align}
với các giá trị $x_k, y_k, k=\overline{1,n}$ là dữ liệu đã có và đóng vai trò hằng số trong hệ phương trình tuyến tính này. Khi giải hệ phương trình \eqref{eq_linear_leastsquare}, ta sẽ tìm được các giá trị $A^*$ và $B^*$ của hàm xấp xỉ $y=A^* + B^*x$.
}

\medskip

Tương tự, đối với dạng hàm $f(x)=A p(x) + B q(x)$, ta lập hàm số $g(A,B) = \sum_{k=1}^n [Ap(x_k)+Bq(x_k)-y_k]^2$, và tính các đạo hàm $\frac{\partial g}{\partial A}, \frac{\partial g}{\partial B}$, để tìm ra giá trị $A^*,B^*$ mà tại đó tổng sai số $g(A^*,B^*)$ là nhỏ nhất.

\newpage\titls{Bài toán xấp xỉ hàm thực nghiệm}

\begin{sagesilent}
a0=2
b0=0.8
n=10
xk=range(n)
yk=[a0+b0*i+random() for i in range(n)]
zk=zip(xk,yk)
plotdata=list_plot(zk, size=50, legend_label='($x_k$,$y_k$)')
xk2=[xk[i]^2 for i in range(n)]
matA=matrix([[n,sum(xk)],[sum(xk),sum(xk2)]])
xkyk=[xk[i]*yk[i] for i in range(n)]
matB=vector([sum(yk),sum(xkyk)])
a,b=N(matA\matB)
plotline=plot(a+b*x,0,n, color='red', legend_label='($y=A+Bx$)')

tableVD5=r"\begin{tabular}{l|c|l}"
tableVD5+=r"$k$ & $x_k$ & $y_k$ \\ \hline"
for i in range(n):
  tableVD5+=latex(i) + r"&" + latex(xk[i]) + r"&" + latex(round(yk[i],3)) + r"\\"
tableVD5+=r"\end{tabular}"
\end{sagesilent}

Ví dụ: lập phương trình đường thẳng $y=A+Bx$ để xấp xỉ tập dữ liệu trong bảng dưới: (phương trình đường thẳng là $y=\sage{round(a,3)}+\sage{round(b,3)}x$)

 \begin{tabular}{lc}
   \sagestr{tableVD5}
  & 
    \begin{tabular}{c}
    \sageplot[scale=0.4]{plotdata+plotline}
    \end{tabular}
 \end{tabular}

%\newpage\titls{Bài tập}
% 
% \newpage\titls{Ví dụ bài tập nhóm}
% 
% \begin{sagesilent}
%  n=7
%  xk3=[0,1,2,3,4,5,6]
%  fxk3=[1,3,5,0,2,6,3]
% # yk3=[(xk3[i],fxk3[i]) for i in range(n)]
%  yk3=zip(xk3,fxk3) # zip: couple pairs of values of the two arrays
%  R = PolynomialRing(QQ, 'x')
%  Lagrange_3 = R.lagrange_polynomial(yk3)
%  plotL3=plot(Lagrange_3, (xk3[0]),(xk3[n-1]), color='red', legend_label='$\mathcal{L}(x)$')
%  plotfxk3=list_plot(yk3, size=50, legend_label='($x_k$,$y_k$)')
%  tableVD3=r"\begin{tabular}{l|c|l}"
%  tableVD3+=r"$k$ & $x_k$ & $y_k$ \\ \hline"
%  for i in range(n):
%    tableVD3+=latex(i) + r"&" + latex(round(xk3[i],2)) + r"&" + latex(fxk3[i]) + r"\\"
%  tableVD3+=r"\end{tabular}"
% \end{sagesilent}
% 
% Ví dụ: tìm hệ số Lagrange của hàm số $f(x)$ có các giá trị cho trong bảng số:
% 
%  \begin{tabular}{lc}
%    \sagestr{tableVD3}
%   & 
%     \begin{tabular}{c}
%     \sageplot[scale=0.3]{plotfxk3+plotL3}
%     \end{tabular}
%  \end{tabular}
% 
% Kết quả: $\sage{Lagrange_3}$
% 
% \newpage\titls{Ví dụ bài tập nhóm}
% 
% \begin{sagesilent}
%  n=7
%  xk3=[0,1,2,3,4,5,6]
%  fxk3=[1,3,5,0,3,1,5]
% # yk3=[(xk3[i],fxk3[i]) for i in range(n)]
%  yk3=zip(xk3,fxk3) # zip: couple pairs of values of the two arrays
%  R = PolynomialRing(QQ, 'x')
%  Lagrange_3 = R.lagrange_polynomial(yk3)
%  plotL3=plot(Lagrange_3, (xk3[0]),(xk3[n-1]), color='red', legend_label='$\mathcal{L}(x)$')
%  plotfxk3=list_plot(yk3, size=50, legend_label='($x_k$,$y_k$)')
%  tableVD3=r"\begin{tabular}{l|c|l}"
%  tableVD3+=r"$k$ & $x_k$ & $y_k$ \\ \hline"
%  for i in range(n):
%    tableVD3+=latex(i) + r"&" + latex(round(xk3[i],2)) + r"&" + latex(fxk3[i]) + r"\\"
%  tableVD3+=r"\end{tabular}"
% \end{sagesilent}
% 
% Ví dụ: tìm hệ số Lagrange của hàm số $f(x)$ có các giá trị cho trong bảng số:
% 
%  \begin{tabular}{lc}
%    \sagestr{tableVD3}
%   & 
%     \begin{tabular}{c}
%     \sageplot[scale=0.3]{plotfxk3+plotL3}
%     \end{tabular}
%  \end{tabular}
% 
%  Kết quả: $\sage{Lagrange_3}$
%  
\label{lastpage}
\end{document}
